\documentclass[a4paper,10pt]{article}
\usepackage[left=3cm,top=3cm,right=2.5cm,bottom=4cm]{geometry}
\usepackage[utf8]{inputenc}
\usepackage[activeacute,spanish]{babel}
\usepackage{multirow}	% Para poder unir filas en las tablas
\usepackage{colortbl}	% Para colorear tablas
\usepackage{latexsym,amssymb,oldlfont,amsmath}
\usepackage{graphicx,graphics}
\usepackage{multicol}
\usepackage{color}
\usepackage{listings}
\usepackage{lastpage}
\usepackage{enumitem}



%\usepackage{fontspec}
%\setmainfont{Times New Roman}
\newcommand{\R}{{\mathbb R}}
\newcommand{\si}{{\mathcal S}^n}
\renewcommand{\theequation}{\thesection.\arabic{equation}} % Cambia formato salida numeracion equation
\makeatletter\@addtoreset {equation}{section} \makeatother % Tambien se usa
\renewcommand{\baselinestretch}{1.3}
\usepackage{amssymb}
\usepackage{hyperref}
\definecolor{midnightblue}{rgb}{0.1,0.1,0.44}


\newcommand{\dis}{\displaystyle}
\renewcommand{\theenumii}{\arabic{enumii}}

\def\Z{\mathbb{Z}}
\def\N{\mathbb{N}}
\def\R{\mathbb{R}}
\def\M{\mathcal{M}_n(\mathbb{R})}
\def\sin{\operatorname{sen}}
\def\e{\operatorname{e}}
\DeclareMathOperator{\arcsec}{arcsec}
\DeclareMathOperator{\arccot}{arccot}
\DeclareMathOperator{\arccsc}{arccsc}



\usepackage{etoolbox,fancyhdr,xcolor}
\fancyhf{} % sets both header and footer to nothing
\renewcommand{\headrulewidth}{0pt}
%\renewcommand{\footrulewidth}{0.5pt}
\pagestyle{fancy}
\headheight=60pt %para cambiar el tamaño del encabezadox
\fancyhead[L] %la "L" indica a la izquierda
{
\begin{minipage}{6cm}
\includegraphics[width=0.9\textwidth]{logo_udp.eps}
\end{minipage}
}

%\fancyhead[C]
%{
%\textbf{Cálculo II - Guía 1}
%}


%\fancyhead[R] %la "L" indica a la izquierda
%{

%}
%\begin{minipage}{4cm}
%\small %tamaño de letra pequeño
%{
%\textsf
%{
%Universidad de Chile\\
%...}}
%\end{minipage}}



%\fancyfoot[C]{}
%{
%\begin{minipage}{10cm}
%\color{midnightblue}
%\begin{center}Casa Matriz: Monjitas 527 Piso 15 - Santiago de Chile\\
%Fono (02)24367255 \end{center}
%\end{minipage}
%}
\fancyfoot[C]{}
%Página \thepage\ de \pageref{LastPage}}% Páginas Pares numeradas a la derecha
%\fancyfoot[LO]{\thepage} % Páginas Impares numeradas a la izquierda

\title{C2_Ay1_Ej1}

\begin{document}

\begin{center}
Cálculo II\\
Ayudantía Nº1 - Ejercicio Nº1\\
Primer Semestre 2017
\end{center}





1. Calcule las siguientes integrales indefinidas mediante el m\'{e}todo de
sustituci\'{o}n.
\begin{enumerate}[label=\emph{\alph*)}]
\item $\dis \int \sec(x) \tan(x)\sqrt{1+\sec(x)}\, dx$
\item $\dis \int \sqrt[3]{(x^3+1)}\, x^5\, dx$
\end{enumerate}


\vspace{5mm}

Solución:

\begin{enumerate}[label=\emph{\alph*)}]
\item $\dis \int \sec(x) \tan(x)\sqrt{1+\sec(x)}\, dx$\\

Usando la siguiente sustitución:
\begin{align*}
u &= 1+\sec(x) \\
du &= \sec(x) \tan(x) dx
\end{align*}

Obtenemos lo siguiente:
\begin{align*}
 \dis \int \sqrt{1+\sec(x)} \sec(x) \tan(x)\, dx &= \dis \int \sqrt{u} \, du \\
      &=  \frac{2*u^{\frac{3}{2}}}{3} + C\\
       &=  \frac{2*{(1+\sec(x))}^{\frac{3}{2}}}{3} + C
\end{align*}

\textcolor{red}{1 punto por desarrollar completa la a)}

\item $\dis \int \sqrt[3]{(x^3+1)}\, x^5\, dx$\\
Usando la siguiente sustitución:
\begin{align*}
u &= x^3+1 \\
du &= 3x^2 dx
\end{align*}

Despejando $x^3$ obtenemos que: $x^3 = u - 1$

\begin{align*}
 \dis \dis \int \sqrt[3]{(x^3+1)}\, x^5\, dx &= \dis \int \sqrt[3]{(x^3+1)}\, x^3* x^2\, dx \\
      &=  \dis \int \sqrt[3]{u}\, (u-1)\, \frac{du}{3} & \\
 \end{align*}
 \textcolor{red}{1 punto por llegar a la primitiva en función de u en la b)}
 \begin{align*}
       &=  \frac{1}{3}\int (u^{\frac{4}{3}} - u^{\frac{1}{3}}) \, du\\
       &=  \frac{1}{3}\int u^{\frac{4}{3}}\, du - \int u^{\frac{1}{3}} \, du\\
 	   &=   \frac{u^{\frac{7}{3}}}{7}  - \frac{u^{\frac{4}{3}}}{4} + C\\
       &=   \frac{{(x^3+1)}^{\frac{7}{3}}}{7}  - \frac{{(x^3+1)}^{\frac{4}{3}}}{4} + C
\end{align*}
\textcolor{red}{1 punto por volver a la variable x}

Otra manera de hacerlo:

Usando la siguiente sustitución:
\begin{align*}
u^3 &= x^3+1 \\
3u^2 \, du &= 3x^2 dx
\end{align*}

Manipulamos un poco la sustitución para reemplazarla de manera más fácil.
\begin{align*}
x^3 &= u^3-1 \\
u^2 \, du &= x^2 dx
\end{align*}

Obtenemos lo siguiente:

\begin{align*}
 \dis \dis \int \sqrt[3]{(x^3+1)}\, x^5\, dx &= \dis \int \sqrt[3]{(x^3+1)}\, x^3* x^2\, dx \\
      &=  \dis \int \sqrt[3]{u^3}\, (u^3-1) u^2 \, du
       \end{align*}
 \textcolor{red}{1 punto por llegar a la primitiva en función de u en la b)}
 \begin{align*}
       &=  \int u(u^3-1) u^2 \, du \\
       &=  \int u^6\, du - \int u^3 \, du\\
       &=  \frac{u^7}{7} - \frac{u^4}{4} + C\\
        &=  \frac{(x^3+1)^{\frac{7}{3}}}{7} - \frac{(x^3+1)^{\frac{4}{3}}}{4} + C
\end{align*}
\textcolor{red}{1 punto por volver a la variable x}

\end{enumerate}




\end{document}
